\section{Partial Orders (Posets)}

A categoria de ordens parciais é fundamental para modelar hierarquias e dependências funcionais. Um objeto nesta categoria define-se por um conjunto finito e uma relação binária que respeita a reflexividade, antissimetria e transitividade.

\subsection{Objetos}
No ambiente MATLAB, a forma mais eficiente de representar estas relações é através de matrizes de adjacência booleanas. Esta estrutura permite-nos validar propriedades algébricas complexas utilizando operações de álgebra matricial.

\begin{lstlisting}[caption={Representação de um Poset (1 <= 2 <= 3)}]
% Matriz de adjacência para uma ordem linear
n = 3;
M = [1 1 1; 
     0 1 1; 
     0 0 1]; 
\end{lstlisting}

\subsection{Morfismos}
Os morfismos nesta categoria são funções monótonas. Um mapeamento $f$ é válido se, para quaisquer elementos $i, j$, a condição $i \le j$ implicar obrigatoriamente $f(i) \le f(j)$.

\begin{lstlisting}[caption={Algoritmo de Verificação de Monotonia}]
% f é um vetor onde o índice é o domínio e o valor é a imagem
f = [1, 2, 2]; 

isMorphism = true;
[I, J] = find(M_domain); 
for k = 1:length(I)
    if ~M_codomain(f(I(k)), f(J(k)))
        isMorphism = false; break;
    end
end
\end{lstlisting}